% Part: lambda-calculus
% Chapter: introduction
% Section: church-rosser

\documentclass[../../../include/open-logic-section]{subfiles}

\begin{document}

\olfileid{lam}{int}{cr}
\olsection{చర్చ్--రోసర్ లక్షణం}

\begin{thm}
\ollabel{thm:church-rosser}
$M$, $N_1$, $N_2$ పదాలై, $M \red N_1$, $M \red N_2$ అనుకుందాం.
అప్పుడు $N_1 \red P$, $N_2 \red P$ అయ్యే పదం $P$ ఒకటి ఉంటుంది.
\end{thm}

\begin{cor}
$M$ను నియత రూపానికి తగ్గించవచ్చని అనుకుందాం. అప్పుడు ఈ నియత రూపం
ఏకైకం.
\end{cor}

\begin{proof}
$M \red N_1$, $M \red N_2$ అయితే, ముందున్న సిద్ధాంతం ప్రకారం $N_1$,
$N_2$ రెండూ~$P$కు తగ్గేలా ఒక పదం~$P$ ఉంటుంది. $N_1$,~$N_2$ రెండూ
నియత రూపంలో ఉంటే, $N_1 \ident P \ident N_2$ అయినప్పుడే ఇది సాధ్యం.
\end{proof}

చివరిగా, రెండు పదాలు $M$,~$N$ ఒకే పదానికి తగ్గితే, అంటే $M \red P$,
$N \red P$ అయ్యే ఏదో $P$ ఉంటే, వాటిని \emph{$\beta$-తుల్యమైనవి}
లేదా కేవలం \emph{తుల్యమైనవి} అంటాం. దీన్ని $M \equal[\beta] N$ అని
రాస్తాం. $\equal[\beta]$ ఒక తుల్యతా సంబంధమనీ, అంతేకాక ప్రతి $M$,
ప్రతి~$N$కు $M \red N$ లేదా $N \red M$ అయితే $M \equal[\beta] N$
అనీ \olref{thm:church-rosser}ను ఉపయోగించి తనిఖీ చేయవచ్చు. (నిజానికి,
ఈ లక్షణం గల \emph{అత్యల్ప} తుల్యతా సంబంధం $\equal[\beta]$ అని
చూపవచ్చు.)

\end{document}
